% 新しい章の雛形。コピーして chapters/<部>/NN-名前.tex に置き、main.tex の \part の下に \input する。
% 環境・マクロの一覧は docs/design/004-tex-textbook.md の「章の書き方」。
% ファイル名が _ で始まるものは、ビルドと検査の対象外（この雛形が図の検査に引っかからないように）。
\chapter{題}\label{ch:名前}
\chapterinfo{★★☆}{約 X 分}{YYYY-MM-DD}
\correspondsto{参考にする文献の節（なければこの行を消す）}

\begin{objectives}
\begin{enumerate}
  \item （〜を導ける）
  \item （〜を説明できる）
  \item （〜をシミュレーションで確かめられる）
\end{enumerate}
\end{objectives}

\paragraph{前提}
先に読む章は \ref{ch:ode} のように鍵で参照する。必要な数学: （微積分、三角関数、…）。

%% ================================================================
\section{導入}\label{sec:名前-intro}

何が不思議か。日常のどんな現象に対応するか。この章の問いを 1 つ立てる。

\begin{intuition}
この章の問い: …？
\end{intuition}

%% ================================================================
\section{設定}\label{sec:名前-setup}

登場する量、座標の取り方、単位（SI）。図があれば図。

%% ================================================================
\section{本論}\label{sec:名前-main}

方程式を立てる → 解く → 結果の意味を読む。\textbf{途中式は省略しない}。数値は計算してから書く。

\begin{defn}{名前}{鍵}
定義を書く。参照は \ref{def:鍵}。
\end{defn}

\begin{equation}
  \diff{y}{x}=\pdiff{f}{x}\label{eq:名前-a}
\end{equation}

%% ================================================================
\section{シミュレーションで確かめる}\label{sec:名前-sim}

\begin{simbox}
\begin{terminal}
uv run physlab run <コマンド名>
\end{terminal}
動画のどこを見るか。
\end{simbox}

% 図はコマンドの figures() が outputs/figures/ に作る。名前に拡張子は付けない。
% \physfig[0.95]{図の名前}{キャプション}{名前-fig}

理論値と実測値の比較は、\textbf{実際に計算した値}を表にする。

%% ================================================================
\section{演習}\label{sec:名前-exercises}

\begin{exercise}[基礎]{名前-a}
問題文。
\end{exercise}

%% ================================================================
\section{まとめ}\label{sec:名前-summary}

\begin{summary}
\begin{itemize}
  \item 要点 1
  \item 要点 2
\end{itemize}
\end{summary}

\section{発展}\label{sec:名前-advanced}

次に学ぶこと、関連する章・題材。

%% ================================================================
\section{演習の解答}\label{sec:名前-solutions}

\begin{solution}{名前-a}
解答。
\end{solution}
